\section{Conclusions}\label{sec:conclusions}

The central conservation problem addressed in this study was how to direct limited field-inspection capacity across a distributed archaeological property when relevant evidence differs in spatial support, temporal resolution and meaning. The Climate-linked Heritage Inspection Prioritisation (CHIP) framework responded to this problem by combining five Landsat epochs, regional climate records, terrain and drainage context, and explicit sensitivity analysis in a traceable relative-priority architecture. Its purpose is not to replace condition assessment. It identifies where field teams may obtain the greatest immediate value from inspection, while preserving the distinction between satellite-observed landscape pressure, environmental susceptibility and verified condition of archaeological fabric. The framework consequently offers a disciplined way to move from heterogeneous screening evidence to an inspection sequence without presenting analytical neighbourhoods as legal boundaries or relative scores as absolute measures of deterioration.

The evidence demonstrates both the reach and the selectivity of the screening results. Common endpoint support covered 386,280 cells, or 95.7371\% of the analytical grid. Within that support, 65,391 cells (16.9284\%) met at least two adverse spectral criteria, whereas 5,493 cells (1.4220\%) met all three. These proportions show that spectral convergence was spatially concentrated rather than ubiquitous. The Open-Meteo record indicated an annual precipitation trend of $-10.3037$ mm year$^{-1}$ and an annual mean-temperature trend of $+0.041753$ degrees C year$^{-1}$ over 1991 to 2025, but the opposing NASA POWER trend directions and weak annual precipitation agreement prevented a product-independent climatic conclusion. Scene-matched antecedent windows also showed that E2024 was not uniformly the wettest Landsat epoch. Climate therefore supplies a necessary regional and temporal context for inspection planning, but not reliable component-scale differentiation or a basis for attributing spectral change to a particular forcing.

At the primary 500 m analytical support, the declared equal-domain score placed the Giri complex first with a relative priority score of 0.591176. That result is useful only when reported with its uncertainty. Agreement between the 500 and 1,000 m rankings was high (Spearman $\rho=0.918455$), while agreement between 250 and 1,000 m was lower ($\rho=0.710784$). Across 50,000 domain-weight simulations, Giri retained median rank 1 but ranged from rank 1 to 9 across the central 95\% interval. Across 2,000 spatial-block resamples, Bhallar had median rank 1 with an interval from 1 to 5, whereas Giri had median rank 2 with an interval from 1 to 3. These overlapping distributions change the practical interpretation from a single winner to an initial high-priority group. Factor ablation reinforces that caution: removing MNDWI produced no overlap with the full top five and moved one component by 13 places, while the strong NDVI-pressure and NDBI-pressure correlation ($\rho=0.971707$) exposed residual redundancy. Inspection programmes should therefore consider priority tiers, domain-specific inspection checklists and results across more than one spatial support, rather than allocating resources from a fixed rank alone.

The geographically separated proxy experiment provides a further methodological check, not condition validation. With no development-test block overlap and a 2.01 km exclusion buffer, the development-selected multilayer perceptron reproduced WorldCover-derived proxy labels with macro-F1 0.830795 and balanced accuracy 0.810261 on 4,188 held-out samples. These results show transferable land-cover information in the spectral features, but they do not establish accuracy for archaeological condition or validate the CHIP score. The immediate contribution of the study is therefore a reproducible, uncertainty-aware basis for sequencing field inspection at Taxila, with particular attention to components that remain prominent across fixed, scale-sensitive and resampled summaries. The next decisive step is a stratified field campaign that records fabric condition separately from surrounding land cover, tests drainage, moisture, erosion and encroachment interpretations, resolves Saraikala's geometry, incorporates official component polygons, and validates climate and hydrological context with local observations. Until that evidence is available, CHIP should be used as an adaptive inspection-planning framework whose priorities are revised as field observations accumulate.
